% Generated draft: requires mathematical and source review.
\documentclass[11pt]{article}
\usepackage[margin=25mm]{geometry}
\usepackage{fontspec}
\setmainfont{Latin Modern Roman}
\usepackage{amsmath,amssymb,mathtools,unicode-math}
\setmathfont{Latin Modern Math}
\usepackage{enumitem,needspace,xurl,hyperref}
\setlength{\parindent}{0pt}
\setlength{\parskip}{5pt}
\setlength{\emergencystretch}{3em}
\newcommand{\N}{\mathbb{N}}
\newcommand{\Z}{\mathbb{Z}}
\newcommand{\Q}{\mathbb{Q}}
\newcommand{\R}{\mathbb{R}}
\newcommand{\C}{\mathbb{C}}
\newcommand{\F}{\mathbb{F}}
\newcommand{\A}{\mathbb{A}}
\newcommand{\PP}{\mathbb{P}}
\newcommand{\E}{\mathbb{E}}
\newcommand{\eps}{\varepsilon}
\newcommand{\dd}{\,\mathrm{d}}
\begin{document}
\section*{UnsolvedMath ID: 701674}
\textbf{Source problem code:} catalogue entry
\subsection{Definitions and mathematical statement}
A group \(G\) is locally finite if every finitely generated subgroup of \(G\) is finite. A section of \(G\) is a quotient \(H/K\), where \(H\leq G\) and \(K\trianglelefteq H\). Thus the excluded-section hypothesis is
\[
\exists\text{ a finite group }F\quad \forall H\leq G\quad \forall K\trianglelefteq H,\qquad H/K\not\cong F.
\]
Here \(F\) may depend on \(G\); this reads the source's phrase ``a particular finite group'' as the existence of an excluded finite section.

A field \(k\) is locally finite if every finitely generated subfield of \(k\) is finite. For a group \(S\), put \(S^{(0)}=S\) and \(S^{(n+1)}=[S^{(n)},S^{(n)}]\). The group \(S\) is solvable if \(S^{(n)}=1\) for some nonnegative integer \(n\).

Relative to a specified notion of definability, the conventional nontrivial meaning of definable simplicity is that every definable normal subgroup \(N\trianglelefteq G\) satisfies \(N=1\) or \(N=G\). The source instead says that there are no proper definable normal subgroups. Taken literally, that wording excludes every nontrivial group, since \(1\) is definable. The questions below use the conventional meaning, subject to confirmation. The source does not specify the language, an ambient structure, or whether parameters are allowed; the intended definability convention remains unresolved.

The supplied background attributes to Hartley a theorem, based on the classification of finite simple groups, asserting that a simple locally finite group satisfying the excluded-section hypothesis is isomorphic to a group of Lie type over a locally finite field. The following are the proposed extensions of that supplied assertion.

\begin{enumerate}
\item[(a)] Let \(G\) be an infinite definably simple locally finite group satisfying the excluded-section hypothesis. Must \(G\) be isomorphic to a group of Lie type over a locally finite field?
\item[(b)] Let \(G\) satisfy all the hypotheses in part (a), and suppose in addition that
\[
\exists d\in\mathbb Z_{\geq 0}\quad \forall S\leq G,\qquad
S\text{ solvable}\ \Longrightarrow\ S^{(d)}=1.
\]
Must \(G\) be isomorphic to a group of Lie type over a locally finite field?
\end{enumerate}
The bound \(d\) is uniform over all solvable subgroups of \(G\), including infinite ones. Part (b) retains the excluded-section hypothesis by interpreting ``the same question'' as referring to part (a). The source does not specify the precise scope of ``group of Lie type''; that terminology is retained pending review.
\subsection{Short English statement}
The problem asks whether an infinite definably simple locally finite group that omits some finite group as a section must be a group of Lie type over a locally finite field. It asks the same question when the derived lengths of all solvable subgroups have a common finite bound. The meaning of definability and the intended Lie type conclusion require clarification. The supplied registry labels the problem open; this is not an independently verified status assessment.
\subsection{Sources}
\begin{itemize}
\item Local UnsolvedMath registry: supplied complete source record, permanent ID \(701674\), exact source code \texttt{catalogue entry}, title ``Kourovka Notebook Problem 20.16.'' The record attributes the problem to A. V. Borovik, with proposed year 2022, and labels it open. Its supplied update timestamp is \texttt{2026-10-16T14:00:00Z}.
\item The supplied background cites \emph{The Kourovka Notebook (Unsolved Problems in Group Theory)}, 20th issue (2022), and the 21st issue edition (2026), \texttt{arXiv:1401.0300v48}, problem 20.16. The collection metadata names E. I. Khukhro and V. D. Mazurov as editors. These bibliographic details are supplied attribution and have not been independently checked.
\item The source statement attributes its motivating theorem to Hartley and says that it is based on the classification of finite simple groups. No separate bibliographic citation for Hartley's theorem is supplied.
\end{itemize}
\end{document}
